% ==============================================================================
% The Grand Daily Tournament — 2026-09-09 Match (Week 2 Day 3)
% Locked template: EB Garamond + Garamond-Math + STKaiti (v4.1)
% Targeted High-Yield Edition: 2 Math Problems + 2 QM Problems
% ==============================================================================
\documentclass[a4paper,11pt]{article}

\usepackage{fontspec}
\usepackage{amsmath}
\usepackage{unicode-math}
\defaultfontfeatures{Ligatures=TeX}
\setmainfont{EBGaramond}[
  Path           = {c:/texlive/2024/texmf-dist/fonts/opentype/public/ebgaramond/},
  Extension      = .otf,
  UprightFont    = *-Regular,
  ItalicFont     = *-Italic,
  BoldFont       = *-SemiBold,
  BoldItalicFont = *-SemiBoldItalic,
  Ligatures      = TeX,
  Numbers        = OldStyle
]
\setsansfont{LibertinusSans}[
  Path        = {c:/texlive/2024/texmf-dist/fonts/opentype/public/libertinus-fonts/},
  Extension   = .otf,
  UprightFont = *-Regular,
  ItalicFont  = *-Italic,
  BoldFont    = *-Bold,
  Scale       = MatchLowercase
]
\setmathfont{Garamond-Math.otf}[
  Path  = {c:/texlive/2024/texmf-dist/fonts/opentype/public/garamond-math/},
  Scale = MatchLowercase
]

\usepackage{xeCJK}
\xeCJKsetup{
  CJKmath      = false,
  PunctStyle   = kaiming,
  CheckSingle  = true,
  AutoFakeBold = true,
  AutoFakeSlant= false
}
\setCJKmainfont{STKaiti}[
  Path        = {C:/Windows/Fonts/},
  UprightFont = STKAITI.TTF,
  BoldFont    = STKAITI.TTF,
  ItalicFont  = STKAITI.TTF,
  Script      = Default
]
\setCJKsansfont{Microsoft YaHei}
\setCJKmonofont{FangSong}

\usepackage[margin=1.5cm,headheight=14pt,headsep=10pt,footskip=18pt]{geometry}
\usepackage{booktabs,enumitem,xcolor,graphicx,tabularx}

\usepackage[breakable,skins,hooks]{tcolorbox}
\usepackage{tikz}
\usetikzlibrary{calc,arrows.meta}
\usepackage{fancyhdr,lastpage}
\usepackage{microtype}

\definecolor{navy}{HTML}{0F172A}
\definecolor{slate}{HTML}{1E293B}
\definecolor{royal}{HTML}{1E40AF}
\definecolor{mist}{HTML}{64748B}
\definecolor{border}{HTML}{E2E8F0}
\definecolor{paper}{HTML}{F8FAFC}
\definecolor{forest}{HTML}{065F46}
\definecolor{amber}{HTML}{D97706}
\definecolor{wine}{HTML}{991B1B}
\definecolor{ice}{HTML}{EFF6FF}

\linespread{1.08}
\setlength{\parindent}{0pt}
\setlength{\parskip}{3pt plus 1pt minus 1pt}
\setlength{\abovedisplayskip}{4pt plus 1pt minus 1pt}
\setlength{\belowdisplayskip}{4pt plus 1pt minus 1pt}

\pagestyle{fancy}\fancyhf{}
\fancyhead[L]{\textcolor{mist}{\footnotesize\sffamily 2027 Theoretical Physics · 604 Calculus \& 811 Quantum Mechanics}}
\fancyhead[R]{\textcolor{slate}{\footnotesize\sffamily\bfseries Daily Tournament · 2026-09-09}}
\fancyfoot[L]{\textcolor{mist}{\footnotesize\itshape ``From Characteristic Recurrence to Hydrogen Degeneracy.''}}
\fancyfoot[R]{\textcolor{mist}{\footnotesize\sffamily Page \thepage\ of \pageref{LastPage}}}
\renewcommand{\headrulewidth}{0.4pt}
\renewcommand{\headrule}{\color{border}\hrule width\headwidth height\headrulewidth \vskip-\headrulewidth}

\newtcolorbox{briefing}[1]{%
  enhanced, blanker, breakable,
  left=3.5mm, right=2.5mm, top=2.5mm, bottom=2.5mm,
  borderline west={2.5pt}{0pt}{royal},
  colback=paper,
  fonttitle=\sffamily\bfseries\color{slate},
  title={#1},
  attach title to upper={\par\vspace{1.5mm}}
}

\newtcolorbox{problem}[2]{%
  enhanced, breakable,
  colback=white, colframe=border, boxrule=0.5pt, arc=1.5mm,
  left=3.5mm, right=3.5mm, top=2.5mm, bottom=2.5mm,
  fonttitle=\sffamily\bfseries\color{slate},
  colbacktitle=paper, coltitle=slate,
  titlerule=0.3pt, toptitle=1.5mm, bottomtitle=1.5mm,
  title={\raisebox{0pt}{\color{royal}\rule{2.5pt}{8.5pt}}\hspace{4pt}#1\hfill{\footnotesize\color{mist}\itshape #2}}
}

\newtcolorbox{solution}[1]{%
  enhanced, breakable,
  colback=white, colframe=border, boxrule=0.4pt, arc=1mm,
  left=3.5mm, right=3.5mm, top=2.5mm, bottom=2.5mm,
  fonttitle=\sffamily\bfseries\footnotesize\color{forest},
  colbacktitle=paper, coltitle=forest,
  titlerule=0.2pt, toptitle=1.2mm, bottomtitle=1.2mm,
  title={\raisebox{0pt}{\color{forest}\rule{2pt}{7pt}}\hspace{3pt}#1}
}

\newtcolorbox{debrief}[1]{%
  enhanced, breakable,
  colback=paper, colframe=slate, boxrule=0.6pt, arc=1.5mm,
  left=4mm, right=4mm, top=3mm, bottom=3mm,
  fonttitle=\sffamily\bfseries\color{white},
  colbacktitle=slate, coltitle=white,
  titlerule=0pt, toptitle=2mm, bottomtitle=2mm,
  title={#1}
}

\newcommand{\checkbox}{\ensuremath{\mdlgwhtsquare}}
\newcommand{\alerttag}[1]{\textcolor{wine}{\sffamily\bfseries [#1]}}

% ==============================================================================
% PAGE 1: COVER PAGE
% ==============================================================================
\begin{document}
\begin{titlepage}
\thispagestyle{empty}
\begin{tikzpicture}[remember picture,overlay]
  \fill[paper] (current page.south west) rectangle (current page.north east);

  \fill[navy] (current page.north west) rectangle ($(current page.north east)+(0,-3.4cm)$);
  \fill[royal] ($(current page.north west)+(0,-3.4cm)$) rectangle ($(current page.north east)+(0,-3.55cm)$);
  \fill[amber] ($(current page.north west)+(0,-3.55cm)$) rectangle ($(current page.north east)+(0,-3.62cm)$);

  \node[anchor=west, text=white] at ($(current page.north west)+(2.2cm,-1.5cm)$) {
    {\fontsize{12}{14}\selectfont\sffamily\bfseries THE GRAND DAILY TOURNAMENT \textperiodcentered\ 2027}
  };
  \node[anchor=west, text=white!50] at ($(current page.north west)+(2.2cm,-2.3cm)$) {
    {\fontsize{8.5}{11}\selectfont\sffamily UCAS \textperiodcentered\ HIAIS \textperiodcentered\ THEORETICAL PHYSICS ADVANCED TRAINING DOSSIER}
  };

  \begin{scope}[shift={($(current page.center)+(0,1.2cm)$)}]
    \draw[line width=0.5pt, border!80, dashed] (0,0) circle (3.3cm);
    \draw[line width=0.7pt, border] (0,0) circle (2.5cm);
    \draw[line width=0.4pt, royal!30] (0,0) circle (1.6cm);
    \draw[line width=1.2pt, slate!60] (-3.8,-2.2) -- (3.8,2.2);
    \draw[line width=1.2pt, slate!60] (-3.8,2.2) -- (3.8,-2.2);
    
    % Axes
    \draw[line width=0.7pt, royal!65, -{Stealth[length=2.5mm]}] (-4.2,0) -- (4.2,0)
      node[right, font=\footnotesize\sffamily\color{mist}] {$r / a_0$};
    \draw[line width=0.7pt, royal!65, -{Stealth[length=2.5mm]}] (0,-3.6) -- (0,3.6)
      node[above, font=\footnotesize\sffamily\color{mist}] {$E_n,\, P(r)$};

    % Hydrogen discrete energy levels En = -E1/n^2 (horizontal dashed lines)
    \draw[line width=1.1pt, navy, dashed] (-3.2, -2.4) -- (3.2, -2.4) node[right, font=\footnotesize\sffamily\color{navy}] {$n=1$};
    \draw[line width=0.9pt, royal, dashed] (-3.2, -0.6) -- (3.2, -0.6) node[right, font=\footnotesize\sffamily\color{royal}] {$n=2$};
    \draw[line width=0.7pt, amber!90!black, dashed] (-3.2, -0.27) -- (3.2, -0.27) node[right, font=\footnotesize\sffamily\color{amber!90!black}] {$n=3$};
    \draw[line width=0.6pt, forest, dashed] (-3.2, 0) -- (3.2, 0) node[right, font=\footnotesize\sffamily\color{forest}] {$E_\infty=0$};

    % Radial probability density P(r) for 1s
    \draw[line width=1.4pt, royal, smooth, samples=120, domain=0:3.8]
      plot ({\x}, {3.5*\x*\x*exp(-1.8*\x)});

    % Vertical drops for rp and <r>
    \draw[line width=0.6pt, forest, dashed] (1.11, 0) -- (1.11, 1.88);
    \draw[line width=0.6pt, amber!90!black, dashed] (1.67, 0) -- (1.67, 1.62);

    \fill[navy] (0,0) circle (3.5pt);
    \fill[royal] (1.11, 1.88) circle (2.2pt);
    \fill[amber] (1.67, 1.62) circle (2.2pt);

    \node[above, font=\footnotesize\sffamily\color{royal}] at (1.11, 1.95) {$r_p = a_0$};
    \node[above right, font=\footnotesize\sffamily\color{amber!90!black}] at (1.67, 1.65) {$\langle r \rangle = 1.5 a_0$};
  \end{scope}

  \node[anchor=center] at ($(current page.center)+(0,-3.4cm)$) {
    \begin{minipage}{0.84\textwidth}
      \centering
      {\fontsize{30}{36}\selectfont\bfseries\color{navy} Daily Tournament}\\[4pt]
      {\fontsize{11}{14}\selectfont\sffamily\bfseries\color{royal}%
        GRAND ARENA OF THEORETICAL PHYSICS \textperiodcentered\ HYDROGEN DEGENERACY \& ALGEBRA}\\[10pt]
      {\color{navy}\rule{0.45\textwidth}{1.2pt}}\\[8pt]
      {\fontsize{10}{13}\selectfont\color{slate}\itshape
        ``604 Calculus \& Algebra: Leibniz Integral Elimination $\times$ Tridiagonal Characteristic Recurrence\\
        811 Quantum Mechanics: Hydrogen Radial Expectation $\times$ Energy Level Degeneracy Summation''\\[3pt]
        {\small\color{mist}\upshape 4-Problem High-Yield Battle \textperiodcentered\ Closed-book \textperiodcentered\ 80\,min Real-time Challenge}}
    \end{minipage}
  };

  \node[anchor=south] at ($(current page.south)+(0,1.5cm)$) {
    \begin{minipage}{0.86\textwidth}
      \begin{tcolorbox}[
        enhanced, colback=white, colframe=border, boxrule=0.6pt, arc=2mm,
        left=4mm, right=4mm, top=3mm, bottom=3mm
      ]
        \renewcommand{\arraystretch}{1.25}
        \sffamily\small
        \begin{tabularx}{\linewidth}{@{}p{2.6cm}X@{}}
          \textbf{Date} & 2026-09-09 \quad\textbar\quad Week 2 (Day 3)\\
          \textbf{Modules} & \textbf{604}: Integral Elimination + Tridiagonal Det \quad\textbar\quad \textbf{811}: Hydrogen Mean Radius + Degeneracy\\
          \textbf{Error Protocol} & 5-Factor: Knowledge\,\textbar\, Modelling\,\textbar\, Logic\,\textbar\, Calculation\,\textbar\, Time\\
          \textbf{Standards} & Target: 100/100 Full Marks \quad\textbar\quad Time Budget: 80 min\\
        \end{tabularx}
      \end{tcolorbox}
    \end{minipage}
  };
\end{tikzpicture}
\end{titlepage}

% ==============================================================================
% PAGE 2: §1 TACTICAL BRIEFING & WARM-UP RETRIEVAL
% ==============================================================================
\clearpage

\begin{center}
  {\LARGE\bfseries\color{navy} Theoretical Physics Training \textperiodcentered\ Daily Tournament}\\[3pt]
  {\small\sffamily\color{mist} 2026-09-09 \quad\textbar\quad Week 2 (Day 3) \quad\textbar\quad Phase: Hydrogen Degeneracy \& Linear Algebra Power Test}\\[-4pt]
  {\color{navy}\rule{\textwidth}{0.8pt}}\\[-11pt]
  {\color{border}\rule{\textwidth}{0.3pt}}
\end{center}

\vspace{-6pt}

\begin{briefing}{\S 1\quad Tactical Briefing \& Warm-Up Retrieval}
\begin{itemize}[leftmargin=1.5em, itemsep=2pt, topsep=1pt]
  \item \textbf{604 今日靶向}：突破变上限积分拆项求导积法则对消全过程，建立反演二阶初值微分方程标准解法；攻克国科大 604 必考 30 分保分板块——$n$ 阶三对角行列式展开与特征递推通项。
  \item \textbf{811 今日靶向}：氢原子物理深度拓展——基态离核距离期望值 $\langle r \rangle = \frac{3}{2}a_0$ 与最概然半径 $r_p = a_0$ 的几何分布差异；定积分验证库仑势维里定理 $\langle T \rangle = -\frac{1}{2}\langle V \rangle$；严格推导主量子数 $n$ 能级简并度求和 $\sum_{l=0}^{n-1}(2l+1)=n^2$ 与自旋 $2n^2$。
  \item \textbf{Spaced Retrieval 间隔提取 (Day +3 滚动回溯)}：\textbf{604 (昨日核心卡点固化)} 变限拆项求导 $\frac{d}{dx}\int_0^x (x-t)f(t)dt = \int_0^x f(t)dt$；\textbf{604 (Day 2)} 二元极值判别式 $\Delta = AC - B^2$；\textbf{811 (Day 2)} 对称有限深方势阱画圆半径 $R = \frac{\sqrt{2mV_0}a}{\hbar}$。
  \item \textbf{考试规则}：闭卷、独立手推、严禁查阅笔记；先在下方草稿区完成 5--8 分钟白纸公式破冰与间隔提取。
\end{itemize}

\vspace{1mm}
\textbf{Warm-Up \& Spaced Retrieval Checklist} (5--8\,min): without notes, write on scratch paper:\\
\checkbox\ \alerttag{Day +3 604} $\frac{d}{dx}\int_0^x (x-t)f(t)dt = \int_0^x f(t)dt$ (Product rule cancellation: $+xf(x) - xf(x) = 0$)\quad
\checkbox\ \alerttag{Day +3 604} Hessian $\Delta = AC - B^2$ ($\Delta>0, A>0 \implies \min$, $\Delta<0 \implies$ saddle) \quad
\checkbox\ \alerttag{Day +3 811} Symmetric finite well circle: $\xi^2+\eta^2 = R^2 = \frac{2mV_0 a^2}{\hbar^2}$ \quad
\checkbox\ Hydrogen mean radius: $\langle r \rangle = \int_0^\infty r P(r)dr = \frac{3}{2}a_0$ \quad
\checkbox\ Virial theorem for Coulomb: $\langle T \rangle = -\frac{1}{2}\langle V \rangle$ \quad
\checkbox\ Degeneracy formula: $\sum_{l=0}^{n-1} (2l+1) = n^2$
\end{briefing}

\vspace{4pt}

\begin{tcolorbox}[
  enhanced, colback=white, colframe=border, boxrule=0.6pt, arc=2mm,
  left=4mm, right=4mm, top=3mm, bottom=3mm,
  fonttitle=\sffamily\bfseries\color{slate},
  title={\raisebox{0pt}{\color{royal}\rule{3pt}{10pt}}\hspace{4pt}Warm-Up Retrieval Workspace \textbar\ 白纸破冰与间隔提取草稿区},
  colbacktitle=paper, coltitle=slate,
  titlerule=0pt, toptitle=2mm, bottomtitle=2mm
]
\textcolor{mist}{\footnotesize\itshape 5--8 分钟闭卷盲写：默写变限拆项求导积法则对消过程、Hessian 判别式，并写出氢原子简并度求和式与平均半径 $\langle r \rangle$ 积分表达式。}
\vspace{9.0cm}
\end{tcolorbox}

% ==============================================================================
% PAGE 3: §2 CORE PROBLEM SET (4 HIGH-YIELD PROBLEMS)
% ==============================================================================
\clearpage

{\Large\sffamily\bfseries\color{navy} \S 2\quad Core Problem Set (604 Math Duo + 811 QM Duo)}\hfill{\small\sffamily\color{mist}\itshape Timed: 80\,min}

\vspace{5pt}

\begin{problem}{M1 \textbar\ 604: Leibniz Integral Equation to 2nd-Order Nonhomogeneous ODE}{Suggested: 15\,min}
设定义在区间 $[0, +\infty)$ 上的连续函数 $f(x)$ 满足积分方程：
\[
  f(x) = \sin x + \int_0^x (x-t)f(t)\,dt.
\]
\begin{enumerate}[leftmargin=1.5em,itemsep=1pt,topsep=1pt]
  \item 将被积表达式拆开，利用 Leibniz 变限求导定理连续求导两次，证明 $f(x)$ 满足二阶常系数线性微分方程 $f''(x) - f(x) = -\sin x$；
  \item 由原积分方程确定初值条件 $f(0)$ 与 $f'(0)$ 的取值；
  \item 求解该二阶常微分方程初值问题，求出 $f(x)$ 的闭式解析解。
\end{enumerate}
\end{problem}

\vspace{3pt}

\begin{problem}{M2 \textbar\ 604: $n$-th Order Tridiagonal Determinant \& Characteristic Recurrence}{Suggested: 15\,min}
计算下列 $n$ 阶三对角行列式 $D_n$（$n \ge 1$）：
\[
  D_n = \begin{vmatrix}
    5 & 2 & 0 & \dots & 0 & 0 \\
    2 & 5 & 2 & \dots & 0 & 0 \\
    0 & 2 & 5 & \dots & 0 & 0 \\
    \vdots & \vdots & \ddots & \ddots & \vdots & \vdots \\
    0 & 0 & 0 & \dots & 5 & 2 \\
    0 & 0 & 0 & \dots & 2 & 5
  \end{vmatrix}_n
\]
\begin{enumerate}[leftmargin=1.5em,itemsep=1pt,topsep=1pt]
  \item 按第一行展开，导出二阶线性齐次递推关系式 $D_n = 5D_{n-1} - 4D_{n-2}$（$n \ge 3$）；
  \item 计算初始值 $D_1$ 与 $D_2$；
  \item 建立特征方程 $r^2 - 5r + 4 = 0$，求解特征根，求出 $D_n$ 关于阶数 $n$ 的通项显式闭式解。
\end{enumerate}
\end{problem}

\vspace{3pt}

\begin{problem}{Q1 \textbar\ 811: Hydrogen Radial Expectation Values \& Virial Theorem}{Suggested: 25\,min}
已知处于基态的氢原子波函数为 $\psi_{100}(r) = \frac{1}{\sqrt{\pi a_0^3}} e^{-r/a_0}$，其中 $a_0 = \frac{\hbar^2}{\mu e^2}$ 为玻尔半径。
\begin{enumerate}[leftmargin=1.5em,itemsep=1pt,topsep=1pt]
  \item 利用径向几率密度 $P(r) = 4\pi r^2 |\psi_{100}|^2$，定积分计算电子离核距离的平均值（期望值）$\langle r \rangle = \int_0^\infty r P(r)\,dr$；
  \item 结合昨日求得的最概然半径 $r_p = a_0$，简述从物理图像与几率分布的长尾几何特征来看，为什么 $\langle r \rangle = \frac{3}{2}a_0 > r_p$；
  \item 定积分计算库仑相互作用势能期望值 $\langle V \rangle = \langle -e^2/r \rangle = \int_0^\infty (-\frac{e^2}{r}) P(r)\,dr$；结合基态总能量 $E_1 = -\frac{e^2}{2a_0}$，求动能期望值 $\langle T \rangle$，并验证维里定理 $\langle T \rangle = -\frac{1}{2}\langle V \rangle$。
\end{enumerate}
\end{problem}

\vspace{3pt}

\begin{problem}{Q2 \textbar\ 811: Hydrogen Spectrum, Degeneracy Summation \& Angular Momentum}{Suggested: 25\,min}
考虑处于球对称库仑引力势 $V(r) = -e^2/r$ 中的氢原子体系。
\begin{enumerate}[leftmargin=1.5em,itemsep=1pt,topsep=1pt]
  \item 给定主量子数为 $n$，写出轨道角动量量子数 $l$ 与磁量子数 $m$ 的可能取值；
  \item 不计电子自旋时，利用等差数列求和严格证明主量子数为 $n$ 的能级简并度 $g_n = \sum_{l=0}^{n-1} (2l+1) = n^2$；若计入电子自旋 $s=1/2$，简并度为何是 $2n^2$？简述库仑势能级与 $l$ 无关的“偶然简并”物理根源；
  \item 设氢原子处于归一化定态叠加态 $\psi(r,\theta,\phi) = \frac{1}{\sqrt{3}}\psi_{200} + \sqrt{\frac{2}{3}}\psi_{211}$：
    \begin{enumerate}[leftmargin=1.2em,label=(\alph*)]
      \item 测量轨道角动量平方 $\hat{L}^2$ 与分量 $\hat{L}_z$，分别可能得到什么确定值？各自对应的几率是多少？
      \item 计算角动量分量期望值 $\langle L_z \rangle$。
    \end{enumerate}
\end{enumerate}
\end{problem}

% ==============================================================================
% PAGE 4: §3 MODEL SOLUTIONS & COMMENTARY
% ==============================================================================
\clearpage

{\Large\sffamily\bfseries\color{navy} \S 3\quad Model Solutions \& Commentary}

\vspace{2pt}

\begin{solution}{M1 \textperiodcentered\ Leibniz Integral Equation with Product Cancellation}
(1) 拆项：$f(x) = \sin x + x\int_0^x f(t)dt - \int_0^x tf(t)dt$。\\
一阶求导（乘积法则）：$f'(x) = \cos x + \left[1\cdot\int_0^x f(t)dt + xf(x)\right] - xf(x) = \cos x + \int_0^x f(t)dt$（前后 $xf(x)$ 必精准对消！）。\\
二阶求导：$f''(x) = -\sin x + f(x) \implies f''(x) - f(x) = -\sin x$。证毕！\\
(2) 初值：在原方程中令 $x=0$ 得 $f(0) = \sin 0 + 0 = 0$；在一阶导数式中令 $x=0$ 得 $f'(0) = \cos 0 + 0 = 1$。\\
(3) 特征方程 $r^2 - 1 = 0 \implies r = \pm 1$，齐次通解 $y_h = C_1 e^x + C_2 e^{-x}$。\\
设特解 $y^* = A\cos x + B\sin x$，代入得 $(-A-A)\cos x + (-B-B)\sin x = -\sin x \implies A=0, B=1/2$。\\
通解 $f(x) = C_1 e^x + C_2 e^{-x} + \frac{1}{2}\sin x$。代入初值：\\
$f(0) = C_1 + C_2 = 0 \implies C_2 = -C_1$；$f'(0) = C_1 - C_2 + 1/2 = 2C_1 + 1/2 = 1 \implies C_1 = 1/4, C_2 = -1/4$。\\
最终解析解：$f(x) = \frac{1}{4}(e^x - e^{-x}) + \frac{1}{2}\sin x = \frac{1}{2}\sinh x + \frac{1}{2}\sin x$。
\end{solution}

\vspace{1.5pt}

\begin{solution}{M2 \textperiodcentered\ $n$-th Order Tridiagonal Determinant \& Characteristic Recurrence}
(1) 按第一行展开：$D_n = 5 M_{11} - 2 M_{12}$。其中 $M_{11} = D_{n-1}$；余子式 $M_{12}$ 第一列为 $(2, 0, \dots, 0)^T$，按该列展开得 $2 D_{n-2}$。\\
因此递推关系为：$D_n = 5D_{n-1} - 4D_{n-2}$（$n \ge 3$）。\\
(2) 初值计算：$D_1 = |5| = 5$；$D_2 = \begin{vmatrix} 5 & 2 \\ 2 & 5 \end{vmatrix} = 25 - 4 = 21$。\\
(3) 特征方程 $r^2 - 5r + 4 = 0 \implies (r-1)(r-4) = 0$，特征根为 $r_1 = 1, r_2 = 4$。\\
通解形式为 $D_n = c_1 \cdot 1^n + c_2 \cdot 4^n = c_1 + c_2 4^n$。代入初始条件：\\
$\begin{cases} n=1: c_1 + 4c_2 = 5 \\ n=2: c_1 + 16c_2 = 21 \end{cases} \implies 12c_2 = 16 \implies c_2 = \frac{4}{3}, \quad c_1 = 5 - 4\cdot\frac{4}{3} = -\frac{1}{3}$。\\
故通项显式闭式解为：$D_n = -\frac{1}{3} + \frac{4}{3}\cdot 4^n = \frac{4^{n+1} - 1}{3}$。（验证：$n=1 \implies \frac{16-1}{3}=5$；$n=2 \implies \frac{64-1}{3}=21$，完全正确！）
\end{solution}

\vspace{1.5pt}

\begin{solution}{Q1 \textperiodcentered\ Hydrogen Radial Expectation Values \& Virial Theorem}
(1) 径向几率密度 $P(r) = 4\pi r^2 |\psi_{100}|^2 = \frac{4}{a_0^3} r^2 e^{-2r/a_0}$。\\
$\langle r \rangle = \int_0^\infty r P(r) dr = \frac{4}{a_0^3} \int_0^\infty r^3 e^{-2r/a_0} dr$。令 $t = \frac{2r}{a_0}$，则 $r = \frac{a_0}{2}t, dr = \frac{a_0}{2}dt$：\\
$\langle r \rangle = \frac{4}{a_0^3} \left(\frac{a_0}{2}\right)^4 \int_0^\infty t^3 e^{-t} dt = \frac{a_0}{4} \cdot 3! = \frac{6a_0}{4} = \frac{3}{2}a_0$。\\
(2) 最概然半径 $r_p = a_0$ 对应几率密度峰值点 $\frac{dP}{dr}=0$；而平均半径 $\langle r \rangle = 1.5 a_0$ 是重心期望值。几率密度分布在 $r \in [0, \infty)$ 上非对称，在 $r > a_0$ 区域具有指数渐近的长尾（Long Tail），将总体几率重心向外侧大幅拉移，故 $\langle r \rangle > r_p$。\\
(3) 势能期望值：$\langle V \rangle = \int_0^\infty (-\frac{e^2}{r}) P(r) dr = -\frac{4e^2}{a_0^3} \int_0^\infty r e^{-2r/a_0} dr$。代入 $t = \frac{2r}{a_0}$ 得：\\
$\langle V \rangle = -\frac{4e^2}{a_0^3} \left(\frac{a_0}{2}\right)^2 \int_0^\infty t e^{-t} dt = -\frac{e^2}{a_0} \cdot 1! = -\frac{e^2}{a_0}$。\\
基态总能量 $E_1 = -\frac{e^2}{2a_0}$。因 $E_1 = \langle T \rangle + \langle V \rangle$，得动能期望值 $\langle T \rangle = E_1 - \langle V \rangle = -\frac{e^2}{2a_0} - (-\frac{e^2}{a_0}) = +\frac{e^2}{2a_0} > 0$。\\
显见 $\langle T \rangle = \frac{e^2}{2a_0} = -\frac{1}{2}\langle V \rangle$，维里定理严格成立！
\end{solution}

\vspace{1.5pt}

\begin{solution}{Q2 \textperiodcentered\ Hydrogen Spectrum, Degeneracy Summation \& Angular Momentum}
(1) 量子数取值：轨道角动量 $l = 0, 1, 2, \dots, n-1$（共 $n$ 个取值）；磁量子数 $m = -l, -l+1, \dots, 0, \dots, +l$（共 $2l+1$ 个取值）。\\
(2) 简并度求和：$g_n = \sum_{l=0}^{n-1} (2l+1) = 2\sum_{l=0}^{n-1}l + \sum_{l=0}^{n-1} 1 = 2 \cdot \frac{(n-1)n}{2} + n = n(n-1) + n = n^2$。\\
计入自旋 $s=1/2$ 时，电子具有自旋向上与向下 2 个自旋态 ($m_s = \pm 1/2$)，总简并度乘 2 即 $2n^2$。\\
库仑势 $V \propto -1/r$ 下能级与 $l$ 无关并非一般球对称中心力场特性，而是库仑场特有的“偶然简并”，本质是由 Runge-Lenz 矢量守恒决定的更高阶 $SO(4)$ 动态对称性。\\
(3) (a) 态矢由 $|2,0,0\rangle$（几率 $|1/\sqrt{3}|^2=1/3$）与 $|2,1,1\rangle$（几率 $|\sqrt{2/3}|^2=2/3$）叠加而成。\\
测量 $\hat{L}^2$：可能测得 $0\hbar^2=0$（几率 $1/3$）或 $1(1+1)\hbar^2=2\hbar^2$（几率 $2/3$）。\\
测量 $\hat{L}_z$：可能测得 $0\hbar=0$（几率 $1/3$）或 $1\hbar=\hbar$（几率 $2/3$）。\\
(b) 期望值：$\langle L_z \rangle = \frac{1}{3}(0\hbar) + \frac{2}{3}(1\hbar) = \frac{2}{3}\hbar$。
\end{solution}

% ==============================================================================
% PAGE 5: §4 DAILY DEBRIEF & REFLECTION
% ==============================================================================
\clearpage

{\Large\sffamily\bfseries\color{navy} \S 4\quad Daily Debrief \& Reflection (Week 2 Day 3 Match)}

\vspace{4pt}

\begin{debrief}{Post-Match Scorecard (fill in before sleep, 5\,min)}
\renewcommand{\arraystretch}{1.2}
\sffamily\small
\begin{tabularx}{\linewidth}{@{}p{2.4cm}X@{}}
  \textbf{604 Calculus} & Time: \underline{\hspace{1.2cm}} min \quad\textbar\quad M1: \checkbox\ PASS \quad M2: \checkbox\ PASS\\
  & Error type: \checkbox\ Product rule cancellation \checkbox\ Determinant expansion \checkbox\ Characteristic eq \checkbox\ Time\\
  \midrule
  \textbf{811 QM} & Time: \underline{\hspace{1.2cm}} min \quad\textbar\quad Q1: \checkbox\ PASS \quad Q2: \checkbox\ PASS\\
  & Error type: \checkbox\ $\langle r \rangle$ integral \checkbox\ Virial sign \checkbox\ Degeneracy sum \checkbox\ $\hat{L}^2, \hat{L}_z$ prob\\
  \midrule
  \textbf{Day +3 Review} & Checklist: \checkbox\ 604 变限求导积法则相消 \quad \checkbox\ 604 Hessian 极值判别 \quad \checkbox\ 811 有限深画圆半径\\
  \midrule
  \textbf{Summary} & \textbf{Total Score:} \underline{\hspace{1.4cm}} / 100 \quad\textbar\quad Total Time: \underline{\hspace{1.4cm}} min\\
  & \textbf{Next Spaced retrieval (Day +3):} 9月12日盲写三对角行列式递推与氢原子平均半径 $\langle r \rangle$ 积分\\
\end{tabularx}
\end{debrief}

\vspace{6pt}

\begin{tcolorbox}[
  enhanced, colback=white, colframe=border, boxrule=0.6pt, arc=2mm,
  left=4mm, right=4mm, top=3mm, bottom=3.5mm,
  fonttitle=\sffamily\bfseries\color{slate},
  title={\raisebox{0pt}{\color{forest}\rule{3pt}{10pt}}\hspace{4pt}Key Derivation Insights \& Traps \textbar\ 突破口与避坑沉淀},
  colbacktitle=paper, coltitle=slate,
  titlerule=0pt, toptitle=1.5mm, bottomtitle=1.5mm
]
\sffamily\small
\textbf{Core Breakthrough (今日核心突破口与关键一步)}:\\[26pt]
\rule{\linewidth}{0.3pt}\\[10pt]
\textbf{Fatal Trap \& Common Error (致命陷阱与避坑警示)}:\\[26pt]
\rule{\linewidth}{0.3pt}\\[10pt]
\textbf{Day +3 Action Item (3 天后间隔重做提取的具体题号与断点)}:\\[20pt]
\rule{\linewidth}{0.3pt}
\end{tcolorbox}

\vspace{4pt}

\begin{center}
  \sffamily\footnotesize\color{mist}
  Theoretical Physics Master Tournament \textperiodcentered\ Keep your handwritten test paper archived in vault.
\end{center}

\end{document}
